In recent years, high-angular-resolution interferometric observations with facilities such as ALMA and NOEMA have revealed streamer-like structures in a growing number of protostellar systems. These asymmetric, filamentary features appear to connect collapsing envelopes with forming disks, and have been proposed as potentially important channels for mass and angular momentum transport during early stages of star formation. Despite their increasing observational prevalence, the physical origin and dynamical role of streamers remain debated, and a unified framework for interpreting their morphology and kinematics has been lacking.

We present a kinematic framework to characterize streamer-like structures in protostellar systems and apply it to four targets with reported asymmetric envelope features (Per-emb-2, Per-emb-50, SCrA, and HLTau) observed with NOEMA and ALMA. For each source, the observational data are placed into a consistent geometric convention by rotating the measured centroids into a common model reference frame using the disk position angle, such that the projected angular-momentum axis aligns with the $+z$ direction. We then extract ordered centroid points $(x_{\rm cen},z_{\rm cen},v_{\rm cen})$ that trace the streamer and provide a noise-suppressed representation of the streamer morphology and kinematics for model comparison.

The model describes streamer material as a density-enhanced flow undergoing gravitationally driven infall with an additional azimuthal component, parameterized by five quantities $(\theta,\phi_0,i,t_{\rm s},\omega)$ that specify the streamer orientation, viewing geometry, collapse timescale, and the rotational strength relative to Keplerian rotation. We constrain these parameters using a hierarchical fitting strategy: (i) a five-dimensional grid search to identify viable regions of parameter space, (ii) an MCMC exploration based on the points method likelihood to efficiently refine the parameter constraints, and (iii) for datasets with full spectral cubes, the global method likelihood that compares the model with the complete PPV structure while remaining computationally efficient. Convergence is evaluated using standard ensemble-sampler diagnostics, and only chains that exhibit stable sampling behavior and well-defined autocorrelation times are retained.

Across all targets, the best-fit solutions demonstrate that a kinematic model combining radial infall and rotational motion can successfully reproduce the large-scale streamer morphology and the global line-of-sight velocity gradients observed in the data. Comparisons between moment-map-based fitting and full cube-based fitting further show that incorporating the complete PPV information yields tighter and more reliable constraints on the streamer geometry and kinematics.

Finally, by inferring characteristic mass accretion rates from the fitted velocity fields, we find that the envelope-to-disk accretion associated with large-scale collapse can exceed the accretion rates previously attributed to individual streamers. This result suggests that streamers trace localized accretion channels within a broader infalling envelope, but are not necessarily the dominant mechanism setting the overall disk mass flux budget. Cases with significant residual discrepancies likely reflect additional physical processes not captured by the present kinematic framework, highlighting the need for further observational and theoretical studies to fully understand the diversity of streamer formation mechanisms.